\section{Discusión}
\label{sec:discusion}

\subsection{Qué significa aquí ``multimodal''}

En la taxonomía de Baltru\v{s}aitis et al.\ \cite{baltrusaitis2019multimodal}, la
arquitectura resultante tiene una lectura precisa: la \emph{representación} y la
\emph{traducción} de modalidades están delegadas a modelos de frontera (Whisper para
voz$\to$texto; ResNet18 para pixeles$\to$latente), mientras que la \emph{alineación}
(pares texto--imagen en el registro conjunto de la EHAM) y la decisión sobre evidencia
(asignación, redirección, rechazo) son operaciones de memoria asociativa.
No hay \emph{fusión} numérica de modalidades: cada modalidad consulta su propio
directorio y el punto de encuentro es la identidad del agente especialista --- una
fusión \emph{a nivel del grupo}, coherente con la teoría de Wegner
\cite{wegner1987transactive}, no a nivel de vectores.

Esta separación es la que hace honesto el resultado de tiempo real: el sistema no es
rápido porque un clasificador atajara el camino, sino porque el enrutamiento
asociativo \emph{ya era} barato (proyección matricial) y lo único costoso ---la
recuperación--- admite asincronía sin pérdida conceptual: el grupo decide en
milisegundos \emph{quién sabe}; el especialista tarda un par de segundos en
\emph{recordar}, igual que en los reportes por lotes.

\subsection{La cadena física de percepción es otro límite de representación}

El reporte de ocho agentes concluyó que los límites del sistema estaban en las
representaciones, no en la maquinaria asociativa.
La prueba física de esta extensión añade un caso más de la misma familia: la cadena
impresora$+$cámara desplaza los momentos de color del estímulo, el latente cae fuera
del soporte del directorio, y el sistema ---correctamente--- rechaza.
Ni el directorio ni el encoder están ``mal'': el estímulo llega de otra distribución.

La mitigación elegida respeta la arquitectura: la renormalización de
momentos~\eqref{eq:colorfix} es una corrección de \emph{frontera} (lentes para el
ojo), no un cambio en las memorias, y su política condicionada ---reintentar solo tras
rechazo--- convierte una limitación medida (el traslape de las distribuciones que
impide todo umbral detector) en una regla sin parámetros que no puede dañar entradas
limpias.
Vale subrayar el patrón metodológico, porque se repite en el proyecto: el fallo
(pantalla de teléfono, colores lavados) primero se \emph{reproduce en simulación}, la
causa se mide, y solo entonces se interviene --- y el rechazo del sistema durante todo
el proceso fue comportamiento correcto, no un defecto a suprimir.

\subsection{La plastilina y el objeto definido}

La caracterización de la Sección~\ref{sec:res_pista} pone números a la imagen
intuitiva del problema de la pista faltante: la memoria hetero-asociativa \emph{sabe}
que hay un objeto (el plano proyectado tiene soporte en todas sus características),
pero lo que tiene es una superposición que debe decodificarse
\cite{morales2025missing}.
El hallazgo más instructivo es el del prototipo emergente: la lectura determinista
``más probable'' columna a columna es \emph{peor} que el promedio de una muestra
aleatoria y muchísimo peor que cualquier método con prueba (7.0--9.0 contra medias
de 0.0--1.2 de ST y 0.0--0.2 de SS sobre 30 corridas).
La razón es estructural, no estadística: el máximo por columna descarta las
correlaciones entre características que el registro conjunto sí almacena, de modo que
el patrón resultante no corresponde a ningún objeto contenido en la memoria.
Recordar, en este modelo, no es leer máximos: es \emph{buscar un patrón conjunto
consistente con la pista}, y la prueba de retro-proyección es el criterio de esa
consistencia.

También merece registro la nota de trazabilidad: el mecanismo de recuperación oficial
del sistema fue siempre sample-and-search --- el mejor método del artículo.
La contribución de esta extensión no fue por tanto ``arreglar'' la recuperación, sino
\emph{exponerla}: implementar RS y ST como alternativas seleccionables, medir las
cuatro lecturas con la métrica común, y hacer visible en ambas interfaces la
distinción entre la masa indeterminada y el objeto definido --- el requisito teórico
que motivó esta parte del trabajo.

\subsection{El laboratorio como instrumento de demostración de Wegner}

Los reportes por lotes \emph{miden} los tres procesos de Wegner; la ventana en tiempo
real los \emph{muestra ocurriendo}.
El directorio de sesión que se llena frase a frase (actualización de directorio), la
asignación visible del grupo al especialista (asignación de información) y la
redirección punto a punto con el TME apagado, incluida la anotación ``se queda: es el
experto'' (coordinación de recuperación), convierten el mecanismo en algo observable
en segundos por una audiencia no técnica.
El costo de esa visibilidad fue de ingeniería de interfaz, no de mecanismo: ninguna
operación de memoria se modificó para el modo en vivo.

\subsection{Limitaciones}

\begin{enumerate}
    \item \textbf{Los números de tiempo real son de un solo equipo.} Las latencias de
    la Tabla~\ref{tab:latencias} dependen del hardware (GPU de gama de entrada para
    Whisper; CPU para el encoder); no se midió en otras configuraciones.
    \item \textbf{El banco físico es pequeño.} La validación de color usa 24 imágenes
    en condiciones simuladas ancladas a una cámara concreta; una validación con
    múltiples cámaras e iluminaciones queda pendiente.
    \item \textbf{La caracterización distribucional cubre cuatro pistas de un
    token.} Las 30 corridas por método dan medias y desviaciones estables, y la
    pista \texttt{animal} ya exhibe la separación ST/SS; pero un banco con pistas
    multi-token y de más dominios daría la curva completa de dificultad.
    \item \textbf{El canal de voz depende de un traductor externo.} Los errores de
    transcripción de Whisper entran al sistema como pistas legítimas; la memoria los
    rechaza si no tienen soporte, pero no puede corregirlos.
    \item \textbf{Herencia de los límites y salvedades del hemisferio visual.} El
    25\% de rechazo residual del experimento base se manifiesta también en vivo; el
    laboratorio no lo empeora, pero tampoco lo resuelve.
    Se heredan igualmente dos salvedades metodológicas del experimento base,
    declaradas en su reporte: la partición de ETH-80 es por imagen y no por objeto
    físico (solape de instancias del 100\%), de modo que el 75.0\% de enrutamiento
    ---y con él las cifras visuales de este laboratorio--- mide generalización a
    \emph{vistas nuevas de objetos conocidos}, no a objetos nuevos; y la tolerancia
    $\xi = 2$ del directorio visual se seleccionó con un barrido medido sobre ese
    mismo conjunto de prueba.
    El estímulo físico de este laboratorio (impresión $+$ cámara) queda de hecho
    \emph{más lejos} de train que el test digital, lo que hace a la demo un
    complemento útil de esa salvedad, no una víctima adicional.
\end{enumerate}
